\documentclass[11pt]{article}
\usepackage[a4paper,margin=25mm]{geometry}
\usepackage{amsmath,amssymb,mathtools}
\usepackage{fontspec,unicode-math}
\setmainfont{Latin Modern Roman}
\setsansfont{Latin Modern Sans}
\setmonofont{Latin Modern Mono}
\setmathfont{Latin Modern Math}
\usepackage{xurl,hyperref}
\hypersetup{colorlinks=true,urlcolor=blue,linkcolor=blue}
\newcommand{\sref}[2]{\hyperlink{source:#1:#2}{[S#2]}}
\newcommand{\cataloguescope}[1]{\paragraph{Mathematical question.}#1}
\begin{document}
% UnsolvedMath ID 8800050; source catalogue code AMR-087-0050
\setcounter{section}{8800049}
\section{Monotonicity of maximal curve point counts in genus}\label{8800050}
{\small\sffamily UnsolvedMath ID 8800050 · Graph Theory}
\subsection{Definitions and mathematical statement}

\paragraph{Shared notation.}$\mathbb N=\{1,2,\ldots\}$, and $\mathbb Z,\mathbb Q,\mathbb R,\mathbb C$ denote the integers, rationals, reals and complex numbers. Any other conventions are specified locally.


For a prime power $q$ and integer $g\ge0$, set $N_q(g)=\max_C\#C(\mathbb F_q)$, where $C$ ranges over smooth projective geometrically connected genus-$g$ curves over $\mathbb F_q$. This is a finite maximum by the Hasse--Weil bound. ``Increasing'' here means nondecreasing; the same source's table allows equal adjacent values, so strict increase is not intended.
\paragraph{Statement recorded in the supplied source.}
For every fixed prime power $q$, is $N_q(g+1)\ge N_q(g)$ for every integer $g\ge0$? \sref{8800050}{2}

\subsection{Short English statement}
Can the largest possible number of rational points decrease when the genus increases by one?
\subsection{Sources}
\begin{description}
\item[\hypertarget{source:8800050:1}{S1}] ULAM.AI and the UnsolvedMath Contributors, UnsolvedMath: A Curated Collection of Open Mathematics Problems, record 8800050 (AMR-087-0050). CC BY 4.0. \url{https://huggingface.co/datasets/ulamai/UnsolvedMath}.
\item[\hypertarget{source:8800050:2}{S2}] Everett W. Howe, \emph{Low-genus curves over finite fields: Problems and variations}, IPAM Number Theory and Cryptography: Open Problems workshop (9--13 October 2006), PDF pages 39--41. \url{https://helper.ipam.ucla.edu/publications/SCWS1/SCWS1_6237.pdf}.
\end{description}
\label{end:8800050}
\end{document}
